\documentclass{article}

\textwidth=6.5in
\textheight=8.5in
\voffset=-54pt
\oddsidemargin=0in
\evensidemargin=0in
\setlength{\parskip}{8pt}
\setlength{\parindent}{0pt}

\usepackage{workbook}

\usepackage[sols]{handout}

\begin{document}

\begin{center}
  \large\textsc{Problem Set 11 - Derivatives of Trigonometric Functions}
  
      \pspace{12pt}
  \begin{problemonly}\large Name: \rule{5cm}{0.15mm}\\ \end{problemonly}
\end{center}

\begin{enumerate}

\item 

\begin{grading}
6 points. 2 per part.
\end{grading}

Use the derivatives of sine and cosine and either the Product Rule +
  Chain Rule or the Quotient Rule to show the following:
  \probonly{}
    \begin{enumerate}
    \item
      \begin{problem}[\vfill]
        $\dfrac{d}{dx} (\sec x) = \sec x \tan x$
      \end{problem}

      \begin{solution}
        We can rewrite $\sec x$ as $\frac{1}{\cos x}$ and use the Quotient
        Rule, or we can rewrite it as $(\cos x)^{-1}$ and use the Chain
        Rule; here, we'll use the latter strategy:
        \begin{align*}
          \frac{d}{dx} (\sec x) &= \frac{d}{dx} \left[ (\cos x)^{-1}
          \right] \\
          &= - (\cos x)^{-2} (-\sin x) \\
          &= \frac{\sin x}{\cos^2 x} \\
          &= \frac{1}{\cos x} \cdot \frac{\sin x}{\cos x} \\
          &= \sec x \tan x
        \end{align*}
      \end{solution}

    \item
      \begin{problem}[\vfill]
        $\dfrac{d}{dx} (\csc x) = -\csc x \cot x$
      \end{problem}

      \begin{solution}
        We can rewrite $\csc x$ as $\frac{1}{\sin x} = (\sin x)^{-1}$ and
        use the Chain Rule:
        \begin{align*}
          \frac{d}{dx} (\csc x) &= \frac{d}{dx} \left[ (\sin x)^{-1}
          \right] \\
          &= - (\sin x)^{-2} (\cos x) \\
          &= - \frac{\cos x}{\sin^2 x} \\
          &= - \frac{1}{\sin x} \cdot \frac{\cos x}{\sin x} \\
          &= - \csc x \cot x
        \end{align*}
      \end{solution}

    \item
      \begin{problem}[\vfill]
        $\dfrac{d}{dx} (\cot x) = - \csc^2 x$
      \end{problem}

      \begin{solution}
        We can rewrite $\cot x = \frac{\cos x}{\sin x}$ and use the
        Quotient Rule:
        \begin{align*}
          \frac{d}{dx} (\cot x) &= \frac{d}{dx} \left( \frac{\cos x}{\sin
              x} \right) \\
          &= \frac{(\sin x)(-\sin x) - (\cos x) (\cos x)}{(\sin x)^2} \\
          &= \frac{ - (\sin^2 x + \cos^2 x)}{\sin^2 x} \\
          &= - \frac{1}{\sin^2 x} \\
          &= - \csc^2 x
        \end{align*}
      \end{solution}

    \end{enumerate}
    \probonly{}

  \probonly{\emph{It's helpful to memorize these so that you can more
      easily tackle derivatives such as the one in 
      {{3(e)}}.}}

\item 

\begin{grading}
5 points.
\end{grading}

  \begin{problem}[\vfill\newpage]
    Use linear approximation of $\tan x$ at $x=0$ in order
    to approximate $\tan 0.1$.
  \end{problem}

  \begin{solution}
    \ideas{The idea of linear approximation is that the line tangent to
      a graph $y = f(x)$ at $x = a$ is a good approximation for the actual
      graph around $x = a$.}
    We first need to pick a function to approximate; the obvious choice is
    $f(x) = \tan x$.  We want to approximate $f(0.1)$, so we should pick a
    value of $a$ near $0.1$ for our approximation; let's use $a = 0$,
    since we can evaluate $f(0)$ exactly.

    We'd like to find the line tangent to the graph of $f(x) = \tan x$ at
    $x = 0$.  The slope of this tangent line is $f'(0)$; since $f'(x) =
    \sec^2 x$, $f'(0) = \sec^2 (0) = \frac{1}{\cos^2 (0)} = 1$.  We also
    know that $(0, f(0)) = (0, 0)$ is a point on the tangent line, so the
    tangent line has equation $y - 0 = 1 (x - 0)$, or $y = x$.  Therefore,
    $\tan x \approx x$ for $x$ near 0.  Plugging in $x = 0.1$ gives $\tan
    (0.1) \approx \boxed{0.1}$.    
  \end{solution}

\item

\begin{grading}
15 points, as follows: 2/2/3/5/3
\end{grading}

Example 21.1 on page 692 in
  Gottlieb may help with the following.

 Differentiate the following; you need not simplify
  your answers.  (\emph{Hint:} You will need logarithmic differentiation
  for one part.)

  \probonly{}
    \begin{enumerate}
    \item
      \begin{problem}[\vfill]
        $y =\cos^2x$
      \end{problem}

      \begin{solution}
        $\displaystyle
        \frac{d}{dx}\left(\fcolorbox{red}{white}{$(\fcolorbox{blue}{white}{$\cos
              x$})^2$}\right)= 2 (\cos x)(-\sin x)$.
      \end{solution}

    \item
      \begin{problem}[\vfill]
        $y = \cos(x^2)$
      \end{problem}

      \begin{solution}
        $\displaystyle \frac{d}{dx}\fcolorbox{red}{white}{$\cos\left(
            \fcolorbox{blue}{white}{$x^{2}$}
          \right)$}=-2x\sin\left(x^{2}\right)$
      \end{solution}

    \item
      \begin{problem}[\stretch{2}]
        $y = x\tan^2 x$
      \end{problem}

      \begin{solution}
        $\displaystyle \frac{d}{dx}
        \left(\fcolorbox{green}{white}{$\displaystyle x$}\cdot
          \fcolorbox{green}{white}{$\displaystyle \tan^2 x$} \right) =
        x\cdot \frac{d}{dx} \left(\fcolorbox{red}{white}{$\displaystyle
            \fcolorbox{blue}{white}{$\displaystyle \tan x$}^ { 2}
            $}\right)+ \tan^2 x = x\cdot 2 \tan x \cdot \sec^2 x + \tan^2
        x$.
      \end{solution}

    \item
      \begin{problem}[\nstretch{3}]
        $y = 7 [\cos(5x)+3]^x$
      \end{problem}

      \begin{solution}
        \ideas{Since the base and exponent both depend on $x$, we should use
          logarithmic differentiation.}  We start with the given equation:
        \begin{align*}
          y &= 7 [\cos(5x)+3]^x \\
          \intertext{Take the natural log of both sides and simplify the
            right side:} %
          \ln y &= \ln \Big( 7 [\cos(5x)+3]^x \Big) \\
          &= \ln 7 + \ln \Big( [\cos(5x)+3]^x \Big) \\
          &= \ln 7 + x \ln \Big( \cos(5x) + 3 \Big) \\
          \intertext{Differentiate both sides with respect to $x$:} %
          \frac{d}{dx} (\ln y) &= \frac{d}{dx} \left[ \ln 7 +
            \fcolorbox{green}{white}{$\displaystyle x$}\cdot
            \fcolorbox{green}{white}{$\displaystyle \ln\Big(\cos (5x)+
              3\Big)$} \right] \\
          \frac{1}{y} \frac{dy}{dx} &= x\cdot \frac{d}{dx}
          \left(\fcolorbox{red}{white}{$\displaystyle
              \ln\left(\fcolorbox{blue}{white}{$\displaystyle \cos (5x)+
                  3$}\right)$}\right)+ \ln\Big(\cos (5x)+ 3\Big) \\
          &= x\cdot \frac{1}{\cos (5x)+ 3}\cdot \left(-\sin (5x)\cdot
            5\right)+ \ln\Big(\cos (5x)+ 3\Big) \\
          &= - \frac{5 x \sin(5x)}{\cos (5x) + 3} + \ln \Big( \cos(5x) + 3
          \Big) \\
          \intertext{Solve for $\frac{dy}{dx}$:} %
          \frac{dy}{dx} &= y \left[ - \frac{5 x \sin(5x)}{\cos (5x) + 3} +
            \ln \Big( \cos(5x) + 3
            \Big) \right] \\
          &= \boxed{7 [\cos(5x)+3]^x \left[ - \frac{5 x \sin(5x)}{\cos
                (5x) + 3} + \ln \Big( \cos(5x) + 3 \Big) \right]}
        \end{align*}        
      \end{solution}

    \item
      \begin{problem}[\stretch{3}]
        $y = \sqrt{\sec(2x^3)}$
      \end{problem}

      \begin{solution}
        \begin{align*}
          \frac{d}{dx} \left(\fcolorbox{red}{white}{$\displaystyle
              \fcolorbox{blue}{white}{$\displaystyle  \sec \left(2 x^ { 3}
                \right)$}^ { \frac{1}{2}} $}\right) = {} & \frac{1}{2}
          \sec \left(2 x^ { 3} \right)^ { -\frac{1}{2}} \cdot \frac{d}{dx}
          \left(\fcolorbox{red}{white}{$\displaystyle \sec
              \left(\fcolorbox{blue}{white}{$\displaystyle 2 x^ { 3}
                  $}\right)$}\right) \\ 
          = {} & \frac{1}{2} \sec \left(2 x^ { 3} \right)^ { -\frac{1}{2}}
          \cdot \sec \left(2 x^ { 3} \right) \tan \left(2 x^ { 3}
          \right)\cdot 6 x^ { 2}  
        \end{align*}
      \end{solution}

    \end{enumerate}

    \probonly{}

\item 

\begin{grading}
10 points. 5 for each part. In both parts, students should either reference the connection with the limit definition of the derivative at $x=0$, or mention that the limits imply that $\sin(x)\approx x$ and $\cos(x)\approx 1$ near $x=0$.
\end{grading}

\begin{enumerate}

\item

\begin{problem}[\stretch{2}]
Find the equation of the tangent line to $y=\sin x$ at $x=0$. What does the limit $\displaystyle \lim_{x\to0} \frac{\sin x}{x} = 1$ have to do with the tangent line?
\end{problem}

\begin{solution}
We start by finding $\frac{dy}{dx}$ when $x=0$. $\frac{dy}{dx}=
\frac{d}{dx}\sin x = \cos x$, so the slope of the tangent line at
$x=0$ is $1$. Using the point $(0,0)$, we find that 
the equation of the tangent line is $\boxed{y=x}$.

We have shown in class that $\displaystyle \lim_{x\to0} \frac{\sin x}{x} = 1$. This is the limit definition of the derivative applied to $\sin x$ at $x=0$, so it gives us the slope of the tangent line at $x=0$. 

This means than for $x$ very small, i.e. very close to 0, $\dfrac{\sin x}{x}
\approx 1$. That is, $\sin x \approx x$. Notice that this \textbf{is} the linear
approximation of $\sin x$ at $x=0$. 
\end{solution}



\item
\begin{problem}[\nstretch{2}]
Find the equation of the tangent line to $\cos x$ at $x=0$.
(You can find this by observation without doing any calculus.) What does the limit $\displaystyle\lim_{x\to0}\frac{\cos x-1}{x}=0$ have to do with the tangent line?
\end{problem}

\begin{solution}
We can see that the tangent line to $y=\cos x$ at $x=0$ is horizontal.
Thus, the equation of the tangent line is $y=1$. 

$\lim\limits_{x\to0}\dfrac{\cos x -1}{x}$ is the definition of the derivative of $\cos x$ at $x=0$, so the slope of the tangent line at $x=0$ is $0$. 

Also, for small $x$, $\dfrac{\cos x -1}{x}
\approx0$ implies $\cos x - 1 \approx 0$, or $\cos x \approx 1$.  
\end{solution}

\end{enumerate}

\item 
% 21.2.16
% Gottlieb 21.2.16

\begin{grading}
5 points.
\end{grading}

  \begin{problem}[\vfill]
    Radians are the best angle unit for doing
    calculus with trig functions, because otherwise the derivatives are more complicated, as we'll see in this problem. Suppose we define a
    function $\operatorname{sindeg}(x)$ that gives the sine of the \textbf{number of degrees} $x$. In terms of our usual function $\sin(\theta)$ that uses radians,
    \[\operatorname{sindeg}(x) =  \sin (\tfrac{\pi
        }{180}x)\] where the factor $\frac{\pi}{180}$ is needed to convert from degrees to radians.\footnote{For example, $\operatorname{sindeg}(30) = \sin
      (\frac{\pi}{6})$.}  Find the derivative of $\operatorname{sindeg}(x)$.
    Is $\operatorname{cosdeg}(x)=\cos(\frac{\pi}{180}x)$ the derivative of $\operatorname{sindeg}(x)$?
  \end{problem}

  \begin{solution} % solution written by ??
    \begin{align*}
    \frac{d}{dx} \operatorname{sindeg}(x)
    &= \frac{d}{dx}  \sin( \tfrac{\pi 
    }{180}x) \\[4pt]
    &= \cos ( \tfrac{\pi}{180}x) \cdot
    \tfrac{\pi}{180} \\[4pt]
    &=\operatorname{cosdeg}(x) \cdot
    \tfrac{\pi}{180}
    \end{align*}
    So, the derivative of $\operatorname{sindeg}(x)$ is not $\operatorname{cosdeg}(x)$.
  \end{solution}  

\item 

\begin{grading}
10 points. 6 for part (a), 4 for part (b).
\end{grading}

In this problem, you'll look at the function $f(x) =
  e^{-x} \cos x$.
  \begin{enumerate}
  \item
    \begin{problem}[\nstretch{3}]
      Find all critical points of $f(x) = e^{-x} \cos x$ on $[0, 2\pi]$. 
    \end{problem}

    \begin{solution}
      First, let's find $f'(x)$; using the Product Rule (and the Chain
      Rule on $e^{-x}$), $f'(x) = - e^{-x} \cos x + e^{-x} (-\sin x) = -
      e^{-x} (\cos x + \sin x)$.  This is never undefined, so we just need
      to figure out when it's 0.  $e^{-x}$ is never 0, so we're solving
      \begin{align*}
        \cos x + \sin x &= 0 \\
        \shortintertext{Let's first solve for $\tan x$:}
        \sin x &= - \cos x \\
        \frac{\sin x}{\cos x} &= -1 \\
        \tan x &= - 1 \\
        \intertext{The solutions of this equation on $[0, 2 \pi]$ are}
        x &= \frac{3 \pi}{4}, \frac{7 \pi}{4}
      \end{align*}
      So the critical points are $\boxed{\frac{3
          \pi}{4}, \frac{7 \pi}{4}}$.
    \end{solution}

  \item
    \begin{problem}[\vfill]
      Find the absolute minimum and maximum values of $f(x) = e^{-x} \cos
      x$ on $[0, 2 \pi]$.  Where do these occur?
      \probonly{\emph{Check by graphing $f(x)$ using a calculator or computer.}} 
    \end{problem}

    \begin{solution}
      The easiest strategy is to evaluate $f$ at the critical points and the endpoints.
      \begin{align*}
        f(0) &= 1 \\
        f\left(\frac{3\pi}{4}\right) &=
        -\frac{1}{\sqrt2}e^{-3\pi / 4} \\
        f\left(\frac{7\pi}{4}\right) &=
        \frac{1}{\sqrt2}e^{- 7\pi / 4}\\
        f(2\pi) &= e^{-2\pi}
      \end{align*}
      The smallest of these values is $f \left( \frac{3 \pi}{4} \right)$,
      as that's the only negative value.  The largest of these values is
      $1$ (it's clear the other three values are $< 1$).  
      \begin{framed}
        So, the absolute maximum value is $1$ (occurring at $x = 0$), and
        the absolute minimum value is $-\frac{1}{\sqrt2}e^{-3\pi / 4}$
        (occurring at $x = \frac{3 \pi}{4}$).
      \end{framed}
    \end{solution}

  \end{enumerate}
\end{enumerate}


\psreflection{
\textbf{Reflect Back.} Your first midterm exam is on Tuesday, March 11th,
    in a little over two weeks. You should reflect back on what you've learned and do
    at least one practice exam.

    How to study for an exam that covers several weeks of material:
    \begin{itemize}[topsep=0pt]
        \item Think about what we've studied and how the material
        is connected. How would you summarize what we've learned to someone who hasn't taken the class? Look for themes and problem-solving strategies.

        \item The best way to study is to do problems: that's why
        you've got several practice exams. You can use one of them as
        a diagnostic to guide your studying. Do \textbf{not} rely on notes or solutions while doing problems. This will give you a false sense of your ability. They are learning resources, not problem-solving aids.
    \end{itemize}
    More guidelines can be found on the exams page. Start studying early.

For next time, read \S 21.3 in Gottlieb.
}
  
\end{document}